% Round 20. RT20. family20. New candidate; does not alter the frozen RB18 or DU20 sources.
\section{Ramified toroidal strata and actual Chow-motive closure}
\label{r20rt:section}

\paragraph{Scope and relation to the frozen results.}
RB18's accepted height theorem has an explicit fourfold positivity
scope.  We do not change that scope.  We reconstruct its
dimension-independent stratum geometry here and apply it instead
to actual Chow-motive presentations.  The additional algebraic
point is that the property for \emph{all} closed strata survives
both directions of the specified smooth toroidal modifications.
This gives a separate ramified closure theorem and, after review,
an extension of DU20 to potentially-good coefficients.

All residue-field varieties below are native smooth projective
varieties.  A closed stratum means a connected component of an
intersection of boundary components, including the ambient variety
when the intersection is empty.  Let $\mathcal C(k)$ denote
either the category of split Tate rational Chow motives over $k$,
or the rational Chow motives of abelian type as defined in AT20.
Both are closed under finite direct sums, direct summands,
products and Tate twists.  Every identity below is in actual
rational Chow theory, unless a cohomological realization is
explicitly mentioned.

\subsection{The two directions of boundary-stratum blowup}

\begin{lemma}[Exact reversible all-strata motive criterion]
\label{r20rt:reversible}
Let $(W,\partial W=\bigcup_iD_i)$ be a smooth projective pair
with strict normal-crossings boundary over a field $k$.
Blow up a smooth connected component $Z$ of $D_K=\bigcap_{k\in K}D_k$
with $K\ne\varnothing$, so the center is a proper closed boundary
stratum, with its full reduced transformed boundary.  Every old closed
stratum has motive in $\mathcal C(k)$ if and only if every new
closed stratum does.
\end{lemma}

\begin{proof}
It suffices to work with the connected strata meeting $Z$;
the other strata are unchanged.  Write $c=|K|$, $F$ for
the exceptional boundary component and $\widetilde D_i$ for
strict transforms.  In the convention that $\mathbf P^{\rm lines}$
parametrizes one-dimensional subspaces, the actual normal bundle is
\[
 N_{Z/W}=\bigoplus_{k\in K}\mathcal O_W(D_k)|_Z.
 \tag{RT20.1}\label{r20rt:normal}
\]
This is the dual of the conormal splitting of the transverse
Cartier equations.  The Rees charts used in AT20 give, for
$K\not\subset J$,
\[
 \bigcap_{j\in J}\widetilde D_j
   =\operatorname{Bl}_{D_{J\cup K}}D_J,
 \qquad
 \operatorname{codim}_{D_J}D_{J\cup K}=|K\setminus J|,
 \tag{RT20.2}\label{r20rt:strict}
\]
on the relevant connected components.  For $K\subset J$
this intersection is empty.  The exceptional intersections are
\[
 F\cap\bigcap_{j\in J}\widetilde D_j
  =\mathbf P^{\rm lines}_{D_{J\cup K}}
       \left(\bigoplus_{k\in K\setminus J}
                     \mathcal O_W(D_k)|_{D_{J\cup K}}\right),
 \tag{RT20.3}\label{r20rt:exceptional}
\]
with rank zero interpreted as empty.  An empty center in
\eqref{r20rt:strict} means the unchanged stratum.  The
Cartier-center case is allowed and just renames its boundary.

For the forward implication, the blowup isomorphism is
induced in Chow groups by
\[
 (\alpha,(\beta_a))\longmapsto
 f^*\alpha+\sum_{a=1}^{c-1}j_*(\xi^{a-1}p^*\beta_a),
 \tag{RT20.4}\label{r20rt:chowmap}
\]
and the projective-bundle isomorphism by
$ (\beta_a)\mapsto\sum_a\xi^a p^*\beta_a$.
These formulas after every smooth projective test product
give the corresponding actual Chow-motive decompositions.
Thus \eqref{r20rt:strict} and \eqref{r20rt:exceptional}
put every new closed stratum in $\mathcal C(k)$.

For the reverse implication, fix an old connected stratum
of $D_J$.  If $K\not\subset J$, its strict transform is
the smooth blowup in \eqref{r20rt:strict}.  The original
stratum is an actual Chow summand of that transform: the
graph pullback and graph pushforward compose to the
identity because the proper birational map has degree one.
This includes $J=\varnothing$, hence the ambient $W$.

If $K\subset J$, choose one $k_0\in K$ and put
$J'=J\setminus\{k_0\}$.  Formula
\eqref{r20rt:exceptional} becomes the rank-one bundle
\[
 F\cap\bigcap_{j\in J'}\widetilde D_j
     =\mathbf P^{\rm lines}_{D_J}
                    (\mathcal O_W(D_{k_0})|_{D_J})
     \simeq D_J.
 \tag{RT20.5}\label{r20rt:recovery}
\]
It is an actual new closed stratum, so gives the required
old motive directly.  Strata in other connected components
of $D_K$ were not modified and are treated as unchanged.
This proves the reverse implication without extracting a
summand from a numerical or cohomological motive.
\end{proof}

\subsection{The actual ramified model and its strata in all dimensions}

Let $\mathcal X/R$ be regular projective reduced strictly
semistable of relative dimension $n$, with finite residue field.
Write its components $D_1,\ldots,D_m$ and
$I_i=\mathcal O_{\mathcal X}(-D_i)$.  Let $L/K$ be finite
of ramification index $e$, with uniformizer $t$.
Saito's Lemmas 1.2.4--1.2.5 state that the \emph{actual blowup}
of $\mathcal X_{R_L}$ in
\[
 \mathfrak a_e=
 \prod_{0\le j<k\le m}\prod_{h=1}^{e-1}
       (I_{j+1}\cdots I_k\mathcal O_{\mathcal X_{R_L}}
                                          +t^h\mathcal O_{\mathcal X_{R_L}})
 \tag{RT20.6}\label{r20rt:saito}
\]
is reduced strictly semistable.  Its generic fiber is $X_L$,
and it is projective.  There is no tameness assumption.
For $e=1$ the ideal is the unit ideal.  We denote this model
by $\mathcal Y/R_L$.

\begin{lemma}[Dimension-independent actual toroidal comparison]
\label{r20rt:comparison}
After a finite constant-field splitting, let $S$ be a connected
closed stratum of $\mathcal Y_k$ defined by $q$ components.
It is connected, as a smooth projective strict toroidal pair,
by smooth projective blowups and blowdowns along closed
boundary strata to a product of $b$ projective-line bundles
over an old closed stratum $D_I$, where
$s=|I|\geq q$ and $b=s-q$.
The old-stratum map is native and birational over the common
normal torus torsor; no finite Kummer cover of $D_I$ occurs.
\end{lemma}

\begin{proof}
We include the geometric argument to make its dimension scope
explicit, rather than extending RB18's fourfold conclusion by
its name.  Split the finitely many native connected strata
by a constant extension if necessary and keep all boundary
labels.  A further finite constant extension can be taken
for the finitely many new strata as well.

\emph{Step 1: primitive local orbit lattices.}
Near an old $s$-fold stratum the actual etale chart is
$x_1\cdots x_s=\pi$.  After base change, write
$\pi=ut^e$ and absorb the unit $u$ into one $x_i$; no
root of $u$ is extracted.  The character lattice and its dual are
\[
 M_I=(\mathbf Z^I\oplus\mathbf Z[t])/
                    \mathbf Z((1)_{i\in I},-e),
 \qquad
 N_I=\{(a,h):\sum_i a_i=eh\}.
 \tag{RT20.7}\label{r20rt:lattice}
\]
Saito's subdivision has height-one rays $(a,1)$ and
unimodular cones.  Let the $q$-dimensional cone $\sigma$
define $S$ generically, and let $I$ be the union of the
supports of its rays.  Every $x_i$ for $i\in I$ contains
a positive exponent of a component defining $S$, so it
vanishes on $S$.  Hence the map factors scheme-theoretically
through $D_I$ and dominates it.
The dense orbit character lattice is
\[
 M_\sigma=M_I\cap\sigma^\perp,
                    \qquad\operatorname{rank}M_\sigma=s-q=b.
 \tag{RT20.8}\label{r20rt:orbit}
\]
The cone rays are part of an integral basis.  Thus this is
the primitive free orbit lattice, and its local open-orbit
ring is its group ring over the ring of $D_I^\circ$.
There is no nontrivial finite index and no residual finite
field or Kummer extension in this description.  The proof
works when the residue characteristic divides $e$.

\emph{Step 2: the actual torsor and extending line bundles.}
Choose a basis of $M_\sigma$ with fixed integer representatives
$((c_{i\nu}),d_\nu)$.  The corresponding open-orbit
coordinates are the residues of the valuation-zero monomials
$t^{d_\nu}\prod_i x_i^{c_{i\nu}}$.  Negative exponents
are used only for invertible orbit characters.  If local
defining equations change by $x_i\mapsto g_i x_i$, the
coordinate is multiplied by
$\prod_i(g_i|_{D_I})^{c_{i\nu}}$.  Thus the open is the
product of the nonzero spaces of the actual line bundles
\[
 \mathcal L_\nu=
   \bigotimes_{i\in I}
       (\mathcal O_{\mathcal X}(D_i)|_{D_I})^{c_{i\nu}}.
 \tag{RT20.9}\label{r20rt:bundles}
\]
The scheme-theoretic factorization in Step 1 permits the
restriction of every transition unit to $D_I$.  Changing a
chosen representative only changes the extension by an old
boundary line bundle; the same open torsor remains fixed.

With the quotient convention, put
\[
 \mathcal P_S=
  \prod_{\nu=1}^{b}{}_{D_I}
        \mathbf P_{D_I}(\mathcal O\oplus\mathcal L_\nu^\vee).
 \tag{RT20.10}\label{r20rt:bundleproduct}
\]
Its open where the two quotient coordinates are nonzero
is the same product of $\mathcal L_\nu^\times$ over
$D_I^\circ$.  Its boundary consists of the zero and infinity
sections and the old boundary pullbacks.

\emph{Step 3: verify the comparison on deeper boundaries.}
At $D_{I\cup J}$, $I\cap J=\varnothing$, the lattices fit
in the actual exact sequence
\[
 0\longrightarrow\mathbf Z^J\longrightarrow
 M_{I\cup J}\cap\sigma^\perp\longrightarrow
 M_I\cap\sigma^\perp\longrightarrow0.
 \tag{RT20.11}\label{r20rt:deeper}
\]
To verify the kernel, subtract the appropriate multiple of
the defining relation in \eqref{r20rt:lattice}; a vector
mapping to zero is left with exactly its $J$ coordinates,
with no relation among them.  Fixed lifts of an orbit basis
and the old boundary coordinates therefore form an integral
basis.  The comparison uses integral monomials in this basis,
pullbacks of old boundary coordinates, and unchanged smooth
base coordinates.  It is consequently a toroidal birational
map of the actual pairs with one common open, compatible
with all overlaps.  It is not inferred merely from two
unrelated normal-crossings compactifications of the same
function field.  The boundary on $S$ is the restriction of
the other components of the strict SNC fiber; all its
connected components are retained.

\emph{Step 4: the precise factorization input.}
Use the corrected integral-conical-complex weak factorization
of Matsuki's erratum, item 2, with the projectivity construction
of AKMW, Theorem 2.7.1.  First take the normalized monomial
graph and a common coherent regular subdivision.  This is
projective and resolves only the supplied toroidal monoid
charts.  The ensuing polyhedral cobordism and its integral
desingularization produce smooth star subdivisions and their
inverses, with projective support functions.  A smooth star
is realized by blowing up the corresponding smooth closed
boundary stratum.  Every intermediate variety remains smooth
and projective, and every intermediate boundary is the full
strict transformed boundary.

Only this integral-combinatorial part is used in positive
characteristic.  Its lattice subdivisions, monomial charts,
and integral support-function line bundles make sense over
the native field in every characteristic.  No general
characteristic-$p$ weak factorization or resolution theorem,
and no withdrawn toric strong-factorization algorithm, is
invoked.  Strict labels and the finite constant splitting
remove lattice self-identifications.  This proves the required
chain in every dimension.  Steps 1--4 are the geometric
argument appearing before the fourfold specialization in
RB18; none of them used the number four.
\end{proof}

\begin{theorem}[Ramified preservation of all-strata Chow motives]
\label{r20rt:ramified}
Suppose all closed strata of $\mathcal X_k$ have motives in
$\mathcal C$ after one finite constant extension.  For every
finite valued $L/K$, the Saito model \eqref{r20rt:saito},
after a further finite unramified extension if needed, has
every closed stratum in the same motive class.
\end{theorem}

\begin{proof}
Fix any new closed stratum $S$ and its pair comparison from
Lemma~\ref{r20rt:comparison}.  Every closed stratum of
$\mathcal P_S$ is a product of some of its projective-line
bundle factors over an old closed stratum $D_{I\cup J}$;
the other factors have been fixed to zero or infinity
sections.  Choosing both sections of one factor gives an
empty intersection.  Thus every nonempty closed stratum
has motive in $\mathcal C$ by the actual projective-bundle
formula.  Apply Lemma~\ref{r20rt:reversible} successively
along the comparison chain, in its indicated directions.
It proves the all-strata assertion for the pair on $S$,
and in particular for $S$ itself.  There are only finitely
many new strata.  The finitely many constant splittings
can be combined into one finite unramified extension of
$L$.  This proves the theorem for the actual new model.
\end{proof}

\subsection{Consequences for the full Green and height domains}

\begin{theorem}[Potentially-good coefficients for supplied abelian strata]
\label{r20rt:green}
Fix the original valued $K/\mathbf Q_p$.  Let $Y/K$ be
smooth projective geometrically integral with a regular
projective reduced strictly semistable model whose closed
strata have Chow abelian type after finite constant extension.
Assume weight--monodromy in every cohomological degree for
$Y$ at every invertible coefficient prime.  Let $T/K$ be
smooth projective geometrically integral with potentially
good projective reduction.  Then $T\times Y$ has full
strong G in every $1\leq r\leq\dim(T\times Y)$, after
every finite valued extension and for all new homological
rational cycles.
\end{theorem}

\begin{proof}
Fix a new cycle over a finite $L/K$.  Choose a finite
extension $M/L$ containing a field over which $T$ has a
smooth projective model.  Apply Theorem~\ref{r20rt:ramified}
to the model of $Y$ over $M$, and then take its finite
unramified splitting extension $M'/M$.  Thus $Y_{M'}$
has an actual projective strict model with the required
component Chow motives, while $T_{M'}$ has a smooth
projective model.

AT20's one-carrier comparison-prime argument and PI20's
numerical generalized inverse give a rational actual
incidence inverse on full geometric cohomology at one
invertible prime.  Tensoring its error with
$\Delta_{T_0}$ proves the full identity on all mixed
classes, without a motive or numerical--homological
comparison assumption on $T_0$.

Weight--monodromy is stable under finite extension: the
monodromy logarithm is multiplied by the nonzero ramification
index, leaving its centered filtration unchanged; on its
graded pieces Frobenius is replaced by the appropriate
residue-degree power, giving the same weight rule over the
new residue field.  Tensoring with the pure good-reduction
cohomology of $T$ retains weight--monodromy.  Hence OW19
kills the actual inertia term, and AK19/CR19 kill the
incidence defect.  VT17 gives G over $M'$ for the specified
pulled cycle, with the low-dimensional endpoint theorems
used if the total dimension is at most two.

Finally use finite actual trace divided by the complete
degree $[M':L]$.  The original absolute value is extended
without renormalization.  This recovers the original
specified support and all actual compact tests over $L$.
The proof began with an arbitrary new cycle over $L$ and
built its model afresh; no old-cycle pullback assumption
or arbitrary proper actual pushforward occurs.
\end{proof}

\begin{corollary}[The additional Drinfeld and mixed-product scope]
\label{r20rt:drinfeld}
DU20 remains valid with its coefficient $T$ assumed only
potentially good projective.  More generally one may take
finitely many Drinfeld quotient factors and finitely many
curve or surface factors having supplied projective strict
models with Chow-abelian-type closed strata after finite
extensions, together with a potentially-good projective
coefficient.  This entire product has full G in every
degree and over every finite extension for every new cycle.
\end{corollary}

\begin{proof}
Pass to the finite etale strict Drinfeld covers used in
DU20.  Their strata are split Tate.  Choose a common finite
field for all supplied models; whenever a further extension
is needed, Theorem~\ref{r20rt:ramified} supplies the new
strict model and preserves its all-strata motive property.
AT20's staircase construction combines the factors, keeping
all closed strata of Chow abelian type.  Ito proves
weight--monodromy for the Drinfeld factors, and the known
curve/surface theorems prove it for the other factors;
Kunneth and the tensor-product monodromy filtration prove
it for their product.  Apply Theorem~\ref{r20rt:green}
and descend the finite etale covers by complete degrees.
\end{proof}

At every finite place of a number-field variety supplied
with such local models, these full new-cycle statements
give the full complementary PH16 height domain, over
all finite number-field extensions, by the same finite
compositum and full-degree descent used in DU20.
This preserves the existing native-good DU20 proof unchanged;
the present corollary is a separately audited extension.

\paragraph{Primary inputs and limits.}
Saito, Lemmas 1.2.4--1.2.5:
\url{https://www.ms.u-tokyo.ac.jp/~t-saito/pp/wmr2r.pdf};
Matsuki's corrected weak factorization, item 2:
\url{https://arxiv.org/pdf/math/0006038};
AKMW, Theorem 2.7.1 and its projective conical construction:
\url{https://personal.math.ubc.ca/~karu/papers/semist.pdf}.
These primary statements and the frozen RB18 source and
independent review were checked.  The motive argument does
not imply invariance of NC17's unrefined compatible quotient,
nor a numerical-kernel cancellation theorem for arbitrary
centers.  No nonreduced initial model, nonsplit lattice
monodromy, general positive-characteristic factorization,
or arbitrary model's Chow motive is covered.  General G
and arbitrary-variety full unconditional heights remain open.
